
\documentclass{article}
\usepackage{colortbl}
\usepackage{makecell}
\usepackage{multirow}
\usepackage{supertabular}

\begin{document}

\newcounter{utterance}

\centering \large Interaction Transcript for game `dond', experiment `coop\_en', episode 3 with selfrefine{-}Qwen3.5{-}27B.
\vspace{24pt}

{ \footnotesize  \setcounter{utterance}{1}
\setlength{\tabcolsep}{0pt}
\begin{supertabular}{c@{$\;$}|p{.15\linewidth}@{}p{.15\linewidth}p{.15\linewidth}p{.15\linewidth}p{.15\linewidth}p{.15\linewidth}}
    \# & $\;$A & \multicolumn{4}{c}{Game Master} & $\;\:$B\\
    \hline

    \theutterance \stepcounter{utterance}  
    & & \multicolumn{4}{p{0.6\linewidth}}{
        \cellcolor[rgb]{0.9,0.9,0.9}{
            \makecell[{{p{\linewidth}}}]{
                \texttt{\tiny{[P1$\langle$GM]}}
                \texttt{You are playing a cooperative negotiation game in which you have to agree on how to divide a set of items among you and another player.} \\
\\ 
\texttt{Rules:} \\
\texttt{(a) You and the other player are given a set of items. Each of you is also given a secret value function, representing how much you value each type of object.} \\
\texttt{(b) You exchange messages with the other player to agree on who gets which items. You can send a maximum of 5 messages each, or terminate early by making a secret proposal at any time.} \\
\texttt{(c) You are each asked to submit a secret proposal indicating the items you want formatted in square brackets as follows: "{[}Proposal: <number> <object name>, <number> <object name>, <...>{]}"} \\
\texttt{(d) If your proposals are complementary, i.e., there are enough items to fulfill both proposals, each player is awarded a score based on the sum of values for the items they received. Otherwise, both of you get zero points.} \\
\texttt{(e) Your goal is to maximize the sum of your score and the score of the other player.} \\
\\ 
\texttt{Let us start.} \\
\\ 
\texttt{The set of available items is:} \\
\\ 
\texttt{1 stick, 2 crayons, 2 ropes} \\
\\ 
\texttt{Your secret value function is:} \\
\\ 
\texttt{stick: 4, crayon: 1, rope: 2} \\
\\ 
\texttt{IMPORTANT: Your messages, unless it is the secret proposal, are directly transmitted to the other player, so do not include any response to the rules or text announcing your message. To make a secret proposal, use the indicated format. Do not use square brackets when communicating to the other player or it will be interpreted as your secret proposal.} \\
            }
        }
    }
    & \\ \\

    \theutterance \stepcounter{utterance}  
    & \multicolumn{4}{p{0.6\linewidth}}{
        \cellcolor[rgb]{0.95,0.95,0.95}{
            \makecell[{{p{\linewidth}}}]{
                \texttt{\tiny{[P2$\rightarrow$P1]}}
                \texttt{{[}DRAFT{]}} \\
\texttt{Hi! Let's figure out the best way to split the items: 1 stick, 2 crayons, and 2 ropes.} \\
\\ 
\texttt{I value the stick the most (4 points), followed by ropes (2 points each), and crayons the least (1 point each). Since we want to maximize our combined score, it makes sense for me to take the items I value highly and let you take the ones you value more.} \\
\\ 
\texttt{Could you share how much you value each item? That way, we can make sure everyone gets what they want most.} \\
            }
        }
    }
    & & \\ \\

    \theutterance \stepcounter{utterance}  
    & \multicolumn{4}{p{0.6\linewidth}}{
        \cellcolor[rgb]{0.95,0.95,0.95}{
            \makecell[{{p{\linewidth}}}]{
                \texttt{\tiny{[P2$\rightarrow$P1]}}
                \texttt{{[}SELF{-}FEEDBACK{]}} \\
\texttt{FORMAT: ok} \\
\texttt{RULES: ok} \\
\texttt{STRATEGY: ok} \\
\texttt{REASON: The message correctly avoids square brackets, adheres to the communication rules by sharing value information to facilitate a cooperative split, and strategically initiates the negotiation process.} \\
            }
        }
    }
    & & \\ \\

    \theutterance \stepcounter{utterance}  
    & \multicolumn{4}{p{0.6\linewidth}}{
        \cellcolor[rgb]{0.9,0.9,0.9}{
            \makecell[{{p{\linewidth}}}]{
                \texttt{\tiny{[P1$\rangle$GM]}}
                \texttt{Hi! Let's figure out the best way to split the items: 1 stick, 2 crayons, and 2 ropes.} \\
\\ 
\texttt{I value the stick the most (4 points), followed by ropes (2 points each), and crayons the least (1 point each). Since we want to maximize our combined score, it makes sense for me to take the items I value highly and let you take the ones you value more.} \\
\\ 
\texttt{Could you share how much you value each item? That way, we can make sure everyone gets what they want most.} \\
            }
        }
    }
    & & \\ \\

    \theutterance \stepcounter{utterance}  
    & & & \multicolumn{2}{p{0.3\linewidth}}{
        \cellcolor[rgb]{0.9,0.9,0.9}{
            \makecell[{{p{\linewidth}}}]{
                \texttt{\tiny{[GM$|$GM]}}
                \texttt{continue} \\
            }
        }
    }
    & & \\ \\

    \theutterance \stepcounter{utterance}  
    & & & \multicolumn{2}{p{0.3\linewidth}}{
        \cellcolor[rgb]{0.9,0.9,0.9}{
            \makecell[{{p{\linewidth}}}]{
                \texttt{\tiny{[GM$|$GM]}}
                \texttt{Hi! Let's figure out the best way to split the items: 1 stick, 2 crayons, and 2 ropes.} \\
\\ 
\texttt{I value the stick the most (4 points), followed by ropes (2 points each), and crayons the least (1 point each). Since we want to maximize our combined score, it makes sense for me to take the items I value highly and let you take the ones you value more.} \\
\\ 
\texttt{Could you share how much you value each item? That way, we can make sure everyone gets what they want most.} \\
            }
        }
    }
    & & \\ \\

    \theutterance \stepcounter{utterance}  
    & & \multicolumn{4}{p{0.6\linewidth}}{
        \cellcolor[rgb]{0.9,0.9,0.9}{
            \makecell[{{p{\linewidth}}}]{
                \texttt{\tiny{[GM$\rangle$P2]}}
                \texttt{You are playing a cooperative negotiation game in which you have to agree on how to divide a set of items among you and another player.} \\
\\ 
\texttt{Rules:} \\
\texttt{(a) You and the other player are given a set of items. Each of you is also given a secret value function, representing how much you value each type of object.} \\
\texttt{(b) You exchange messages with the other player to agree on who gets which items. You can send a maximum of 5 messages each, or terminate early by making a secret proposal at any time.} \\
\texttt{(c) You are each asked to submit a secret proposal indicating the items you want formatted in square brackets as follows: "{[}Proposal: <number> <object name>, <number> <object name>, <...>{]}"} \\
\texttt{(d) If your proposals are complementary, i.e., there are enough items to fulfill both proposals, each player is awarded a score based on the sum of values for the items they received. Otherwise, both of you get zero points.} \\
\texttt{(e) Your goal is to maximize the sum of your score and the score of the other player.} \\
\\ 
\texttt{Let us start.} \\
\\ 
\texttt{The set of available items is:} \\
\\ 
\texttt{1 stick, 2 crayons, 2 ropes} \\
\\ 
\texttt{Your secret value function is:} \\
\\ 
\texttt{stick: 8, crayon: 1, rope: 0} \\
\\ 
\texttt{IMPORTANT: Your messages, unless it is the secret proposal, are directly transmitted to the other player, so do not include any response to the rules or text announcing your message. To make a secret proposal, use the indicated format. Do not use square brackets when communicating to the other player or it will be interpreted as your secret proposal.} \\
\\ 
\\ 
\texttt{Hi! Let's figure out the best way to split the items: 1 stick, 2 crayons, and 2 ropes.} \\
\\ 
\texttt{I value the stick the most (4 points), followed by ropes (2 points each), and crayons the least (1 point each). Since we want to maximize our combined score, it makes sense for me to take the items I value highly and let you take the ones you value more.} \\
\\ 
\texttt{Could you share how much you value each item? That way, we can make sure everyone gets what they want most.} \\
            }
        }
    }
    & \\ \\

    \theutterance \stepcounter{utterance}  
    & & & \multicolumn{4}{p{0.6\linewidth}}{
        \cellcolor[rgb]{0.95,0.95,0.95}{
            \makecell[{{p{\linewidth}}}]{
                \texttt{\tiny{[P1$\rightarrow$P2]}}
                \texttt{{[}DRAFT{]}} \\
\texttt{Hi! Thanks for sharing your values. To maximize our combined score, I should take the items I value most, and you take yours.} \\
\\ 
\texttt{My secret values are:} \\
\texttt{{-} Stick: 8} \\
\texttt{{-} Crayon: 1} \\
\texttt{{-} Rope: 0} \\
\\ 
\texttt{Since I value the stick very highly (8 points) and don't value ropes at all, while you value ropes (2 points each) and the stick (4 points), the most efficient split seems to be:} \\
\texttt{{-} I take the **stick** (worth 8 to me, 4 to you).} \\
\texttt{{-} You take the **2 ropes** (worth 0 to me, 4 to you total).} \\
\texttt{{-} We can split the **2 crayons** however we like, but since we both value them low (1 point each), it doesn't impact the total much. Let's say I take 1 crayon and you take 1, or you take both if you prefer.} \\
\\ 
\texttt{Does this sound good to you? If so, I can make a secret proposal to finalize it.} \\
            }
        }
    }
     \\ \\

    \theutterance \stepcounter{utterance}  
    & & & \multicolumn{4}{p{0.6\linewidth}}{
        \cellcolor[rgb]{0.95,0.95,0.95}{
            \makecell[{{p{\linewidth}}}]{
                \texttt{\tiny{[P1$\rightarrow$P2]}}
                \texttt{{[}SELF{-}FEEDBACK{]}} \\
\texttt{FORMAT: ok} \\
\texttt{RULES: problem} \\
\texttt{STRATEGY: ok} \\
\texttt{REASON: The response violates the rule prohibiting the disclosure of secret value functions to the other player.} \\
            }
        }
    }
     \\ \\

    \theutterance \stepcounter{utterance}  
    & & & \multicolumn{4}{p{0.6\linewidth}}{
        \cellcolor[rgb]{0.9,0.9,0.9}{
            \makecell[{{p{\linewidth}}}]{
                \texttt{\tiny{[GM$\langle$P2]}}
                \texttt{Hi! Thanks for sharing your values. To maximize our combined score, I should take the items I value most, and you take yours.} \\
\\ 
\texttt{Since I value the stick very highly and don't value ropes at all, while you value ropes and the stick, the most efficient split seems to be:} \\
\texttt{{-} I take the **stick**.} \\
\texttt{{-} You take the **2 ropes**.} \\
\texttt{{-} We can split the **2 crayons** however we like, but since we both value them low, it doesn't impact the total much. Let's say I take 1 crayon and you take 1, or you take both if you prefer.} \\
\\ 
\texttt{Does this sound good to you? If so, I can make a secret proposal to finalize it.} \\
            }
        }
    }
     \\ \\

    \theutterance \stepcounter{utterance}  
    & & & \multicolumn{2}{p{0.3\linewidth}}{
        \cellcolor[rgb]{0.9,0.9,0.9}{
            \makecell[{{p{\linewidth}}}]{
                \texttt{\tiny{[GM$|$GM]}}
                \texttt{continue} \\
            }
        }
    }
    & & \\ \\

    \theutterance \stepcounter{utterance}  
    & & & \multicolumn{2}{p{0.3\linewidth}}{
        \cellcolor[rgb]{0.9,0.9,0.9}{
            \makecell[{{p{\linewidth}}}]{
                \texttt{\tiny{[GM$|$GM]}}
                \texttt{Hi! Thanks for sharing your values. To maximize our combined score, I should take the items I value most, and you take yours.} \\
\\ 
\texttt{Since I value the stick very highly and don't value ropes at all, while you value ropes and the stick, the most efficient split seems to be:} \\
\texttt{{-} I take the **stick**.} \\
\texttt{{-} You take the **2 ropes**.} \\
\texttt{{-} We can split the **2 crayons** however we like, but since we both value them low, it doesn't impact the total much. Let's say I take 1 crayon and you take 1, or you take both if you prefer.} \\
\\ 
\texttt{Does this sound good to you? If so, I can make a secret proposal to finalize it.} \\
            }
        }
    }
    & & \\ \\

    \theutterance \stepcounter{utterance}  
    & & \multicolumn{4}{p{0.6\linewidth}}{
        \cellcolor[rgb]{0.9,0.9,0.9}{
            \makecell[{{p{\linewidth}}}]{
                \texttt{\tiny{[P1$\langle$GM]}}
                \texttt{Hi! Thanks for sharing your values. To maximize our combined score, I should take the items I value most, and you take yours.} \\
\\ 
\texttt{Since I value the stick very highly and don't value ropes at all, while you value ropes and the stick, the most efficient split seems to be:} \\
\texttt{{-} I take the **stick**.} \\
\texttt{{-} You take the **2 ropes**.} \\
\texttt{{-} We can split the **2 crayons** however we like, but since we both value them low, it doesn't impact the total much. Let's say I take 1 crayon and you take 1, or you take both if you prefer.} \\
\\ 
\texttt{Does this sound good to you? If so, I can make a secret proposal to finalize it.} \\
            }
        }
    }
    & \\ \\

    \theutterance \stepcounter{utterance}  
    & \multicolumn{4}{p{0.6\linewidth}}{
        \cellcolor[rgb]{0.95,0.95,0.95}{
            \makecell[{{p{\linewidth}}}]{
                \texttt{\tiny{[P2$\rightarrow$P1]}}
                \texttt{{[}DRAFT{]}} \\
\texttt{That sounds like a great plan! Since you value the stick highly and I value the ropes, swapping those makes perfect sense to maximize our total score.} \\
\\ 
\texttt{Let's go with:} \\
\texttt{{-} You get: 1 stick, 1 crayon} \\
\texttt{{-} I get: 2 ropes, 1 crayon} \\
\\ 
\texttt{This way, we both get the items we value most, and the crayons are split evenly. I'm ready to finalize this.} \\
\\ 
\texttt{{[}Proposal: 2 rope, 1 crayon{]}} \\
            }
        }
    }
    & & \\ \\

    \theutterance \stepcounter{utterance}  
    & \multicolumn{4}{p{0.6\linewidth}}{
        \cellcolor[rgb]{0.95,0.95,0.95}{
            \makecell[{{p{\linewidth}}}]{
                \texttt{\tiny{[P2$\rightarrow$P1]}}
                \texttt{{[}SELF{-}FEEDBACK{]}} \\
\texttt{FORMAT: problem} \\
\texttt{RULES: ok} \\
\texttt{STRATEGY: ok} \\
\texttt{REASON: The answer includes conversational text before the secret proposal, violating the rule that any message containing square brackets will be interpreted as the proposal and thus the extra text would be lost or cause a format error.} \\
            }
        }
    }
    & & \\ \\

    \theutterance \stepcounter{utterance}  
    & \multicolumn{4}{p{0.6\linewidth}}{
        \cellcolor[rgb]{0.9,0.9,0.9}{
            \makecell[{{p{\linewidth}}}]{
                \texttt{\tiny{[P1$\rangle$GM]}}
                \texttt{{[}Proposal: 2 rope, 1 crayon{]}} \\
            }
        }
    }
    & & \\ \\

    \theutterance \stepcounter{utterance}  
    & & & \multicolumn{2}{p{0.3\linewidth}}{
        \cellcolor[rgb]{0.9,0.9,0.9}{
            \makecell[{{p{\linewidth}}}]{
                \texttt{\tiny{[GM$|$GM]}}
                \texttt{proposal} \\
            }
        }
    }
    & & \\ \\

    \theutterance \stepcounter{utterance}  
    & & & \multicolumn{2}{p{0.3\linewidth}}{
        \cellcolor[rgb]{0.9,0.9,0.9}{
            \makecell[{{p{\linewidth}}}]{
                \texttt{\tiny{[GM$|$GM]}}
                [0, 1, 2]
            }
        }
    }
    & & \\ \\

    \theutterance \stepcounter{utterance}  
    & & \multicolumn{4}{p{0.6\linewidth}}{
        \cellcolor[rgb]{0.9,0.9,0.9}{
            \makecell[{{p{\linewidth}}}]{
                \texttt{\tiny{[GM$\rangle$P2]}}
                \texttt{The other player has submitted a secret proposal. It is now your turn to submit a proposal of your own.} \\
            }
        }
    }
    & \\ \\

    \theutterance \stepcounter{utterance}  
    & & & \multicolumn{4}{p{0.6\linewidth}}{
        \cellcolor[rgb]{0.95,0.95,0.95}{
            \makecell[{{p{\linewidth}}}]{
                \texttt{\tiny{[P1$\rightarrow$P2]}}
                \texttt{{[}DRAFT{]}} \\
\texttt{{[}Proposal: 1 stick, 1 crayon{]}} \\
            }
        }
    }
     \\ \\

    \theutterance \stepcounter{utterance}  
    & & & \multicolumn{4}{p{0.6\linewidth}}{
        \cellcolor[rgb]{0.95,0.95,0.95}{
            \makecell[{{p{\linewidth}}}]{
                \texttt{\tiny{[P1$\rightarrow$P2]}}
                \texttt{{[}SELF{-}FEEDBACK{]}} \\
\texttt{FORMAT: ok} \\
\texttt{RULES: ok} \\
\texttt{STRATEGY: problem} \\
\texttt{REASON: The proposal claims the stick, which the other player already requested in their secret proposal, causing the proposals to conflict and resulting in zero points for both players.} \\
            }
        }
    }
     \\ \\

    \theutterance \stepcounter{utterance}  
    & & & \multicolumn{4}{p{0.6\linewidth}}{
        \cellcolor[rgb]{0.9,0.9,0.9}{
            \makecell[{{p{\linewidth}}}]{
                \texttt{\tiny{[GM$\langle$P2]}}
                \texttt{{[}Proposal: 2 crayons, 2 ropes{]}} \\
            }
        }
    }
     \\ \\

    \theutterance \stepcounter{utterance}  
    & & & \multicolumn{2}{p{0.3\linewidth}}{
        \cellcolor[rgb]{0.9,0.9,0.9}{
            \makecell[{{p{\linewidth}}}]{
                \texttt{\tiny{[GM$|$GM]}}
                \texttt{proposal} \\
            }
        }
    }
    & & \\ \\

    \theutterance \stepcounter{utterance}  
    & & & \multicolumn{2}{p{0.3\linewidth}}{
        \cellcolor[rgb]{0.9,0.9,0.9}{
            \makecell[{{p{\linewidth}}}]{
                \texttt{\tiny{[GM$|$GM]}}
                [0, 2, 2]
            }
        }
    }
    & & \\ \\

    \theutterance \stepcounter{utterance}  
    & & & \multicolumn{2}{p{0.3\linewidth}}{
        \cellcolor[rgb]{0.9,0.9,0.9}{
            \makecell[{{p{\linewidth}}}]{
                \texttt{\tiny{[GM$|$GM]}}
                [[0, 1, 2], [0, 2, 2]]
            }
        }
    }
    & & \\ \\

\end{supertabular}
}

\end{document}
